\documentclass[a4paper,12pt]{article}

\usepackage[a4paper,top=3.1cm,bottom=3.1cm,left=2.7cm,right=2.7cm]{geometry}
\usepackage[T1]{fontenc}
\usepackage[utf8]{inputenc}
\usepackage{setspace}
\usepackage{graphicx}
\usepackage{booktabs}
\usepackage{array}
\usepackage{float}
\usepackage{amsmath,amssymb,amsthm,mathtools}
\usepackage{enumitem}
\usepackage[round]{natbib}
\usepackage[labelfont=bf,labelsep=period,justification=centering]{caption}
\usepackage{subcaption}
\usepackage{adjustbox}
\usepackage{rotating}
\usepackage{pdflscape}
\usepackage{longtable}
\usepackage{algorithm}
\usepackage{algpseudocode}
\usepackage[colorlinks=true,hyperindex=true,bookmarks=true,pdfpagemode=UseOutlines,unicode]{hyperref}
\hypersetup{citecolor=blue,filecolor=magenta,linkcolor=red,urlcolor=red,
  pdftitle={Internet Appendix: Sequential Completion in Dependent Panels},
  pdfauthor={Giovanni Urga and Emilio Zanetti Chini}}

\DeclareMathOperator*{\argmin}{arg\,min}
\newcommand{\E}{\mathbb{E}}
\newcommand{\Prb}{\mathbb{P}}
\newcommand{\Kcal}{\mathcal{K}}
\newcommand{\Fcal}{\mathcal{F}}
\newcommand{\Ucal}{\mathcal{U}}
\newcommand{\Vcal}{\mathcal{V}}
\newcommand{\Tcal}{\mathcal{T}}
\newcommand{\1}{\mathbf{1}}
\newcommand{\ExhibitNote}[1]{\begin{minipage}{0.94\textwidth}\footnotesize\emph{Notes:} #1\end{minipage}}
\newcommand{\smalltab}{\scriptsize}
\newcolumntype{P}[1]{>{\raggedright\arraybackslash}p{#1}}
% Figure search paths. Keep figures in a local exhibits/ folder for portability.
\graphicspath{{exhibits/}{./exhibits/}{figures/}{./figures/}}

\begin{document}
\pagenumbering{arabic}
\onehalfspacing

\renewcommand{\thetable}{IA.\arabic{table}}
\renewcommand{\thefigure}{IA.\arabic{figure}}
\renewcommand{\thealgorithm}{IA.\arabic{algorithm}}
\newtheorem{iaassumption}{Assumption}
\renewcommand{\theiaassumption}{IA.\arabic{iaassumption}}
\newtheorem{iaproposition}{Proposition}
\renewcommand{\theiaproposition}{IA.\arabic{iaproposition}}
\setcounter{table}{0}
\setcounter{figure}{0}
\setcounter{algorithm}{0}

\begin{center}
{\Large \textbf{INTERNET APPENDIX}}\\[0.5em]
{\large \textit{Sequential Completion in Dependent Panels}}\\[0.5em]
Giovanni Urga and Emilio Zanetti Chini
\end{center}

\addcontentsline{toc}{section}{Internet Appendix}

This Internet Appendix documents the empirical and simulation material that is useful for replication and auditability but not needed for the main line of the paper. The main text states the econometric framework, the formal approximation results, and the headline Monte Carlo and empirical findings. The appendix records the implementation details that a reader would need to reconstruct the exercises: the real-time data layer, the frontier-window construction, the fixed method classes, release-validation diagnostics, target-object timing, bootstrap inference for prespecified empirical comparisons, and the extended simulation output.

All exhibits follow the same convention. The title above each exhibit is intentionally short. The note below the exhibit is longer and states the sample, timing convention, loss object, and interpretation. This keeps the appendix readable while making each table or figure self-contained. The appendix is therefore not a second narrative version of the paper; it is an audit file for the timing, validation, and simulation objects used in the main text.

The appendix is organized as follows. Section~1 documents the real-time FRED-MD data layer and the vintage-index timing convention. Section~2 reports frontier-window diagnostics, showing how the release-validation sample changes when the current-frontier window is widened. Section~3 records the fixed method classes and their real-time feasibility restrictions. Section~4 reports release-validation state diagnostics, state risks, and selector behavior. Section~5 documents the downstream target-object design and the relation between local release loss and INDPRO target-object loss. Section~6 describes bootstrap inference for prespecified empirical comparisons. Section~7 collects the extended Monte Carlo exhibits, including calibration, bridge diagnostics, scaling exercises, application-size regimes, and bridge-strength checks.

\section{Real-Time FRED-MD Data Layer}
\label{ia:data_layer}

The empirical illustration uses a real-time FRED-MD panel organized by data vintages. At vintage $\tau$, the information set $\Fcal_\tau$ contains the releases available at that vintage, the corresponding observation pattern, release-calendar information used in the implementation, and validation losses from earlier decisions whose validation values have already arrived. Candidate completions, states, and selected method classes are therefore vintage-feasible.

A panel cell is indexed by a series $i$ and a reference date $t$. Let $Y_{it}$ denote the validation value fixed by the empirical evaluation design; in this implementation it is the first value observed in a later vintage. Let $r_{it}$ denote the first vintage at which that validation value is available. A cell belongs to the frontier at vintage $\tau$ if it is required for the empirical task and $r_{it}>\tau$. The decision unit is the triple $u=(i,t,\tau)$, because the same calendar cell may be completed at different vintages under different information sets and panel states.

The table below fixes the timing convention used throughout the empirical illustration. The validation date is measured in vintage-index time, not in calendar publication time. This is sufficient for the empirical exercise because validation values, candidate completions, state variables, and feedback availability are all defined relative to the sequence of vintages observed by the econometrician.

\begin{table}[H]
\centering
\caption{Data layer}
\label{tab:ia_data_layer}
\small
\begin{adjustbox}{max width=\textwidth}
\begin{tabular}{@{}ll@{}}
\toprule
Item & Design choice \\
\midrule
Data source & FRED-MD real-time vintages \\
Vintage range & 1999-08--2024-12 \\
Panel date range & 1959-01--2024-11 \\
Number of vintages & 305 \\
Number of panel series & 101 \\
Target variable & INDPRO \\
Decision unit & Frontier triple $u=(i,t,\tau)$ \\
Validation value $Y_{it}$ & First observed later-vintage value \\
Main frontier window & Six months \\
\bottomrule
\end{tabular}
\end{adjustbox}
\par\medskip
\ExhibitNote{The table summarizes the real-time FRED-MD data layer used in the empirical illustration. The validation convention is vintage-index time: a decision unit is validated by the first later vintage in which the prespecified validation value is observed. The main six-month frontier window focuses the empirical exercise on current ragged-edge completion rather than historical backfill-type gaps.}
\end{table}

The arrived validation archive accumulates mechanically as earlier frontier decisions are release-validated in later vintages. In the baseline window used in the empirical illustration, all 4,022 frontier decisions are eventually release-validated and the median reveal delay is one vintage, so the empirical exercise provides frequent delayed full-validation feedback while preserving the decision-date information constraint.

\section{Frontier-Window Diagnostics}
\label{ia:frontier_window}

The main empirical illustration uses a six-month frontier window. This is a design choice, not a tuning outcome. The purpose is to focus on current ragged-edge completion rather than retrospective reconstruction of historical gaps. Wider windows are reported only as diagnostics: they show whether the release-validation sample changes materially when older incomplete cells are admitted.

For a frontier horizon $H$, define
\[
\Ucal^Y_{\tau,H}=\{u=(i,t,\tau):(i,t)\in\mathcal C_\tau,\ r_{it}>\tau,\ 0\le e(\tau)-t\le H\},
\]
where $e(\tau)$ is the last panel reference date at vintage $\tau$. The baseline choice is $H=6$. The diagnostic exhibit compares $H=6$ with wider windows.

\begin{table}[H]
\centering
\caption{Frontier windows}
\label{tab:ia_frontier_windows}
\small
\begin{adjustbox}{max width=\textwidth}
\begin{tabular}{@{}lrrrrrrr@{}}
\toprule
Window & Frontier & Validated & Vintages & Median delay & Mean delay & Never validated & Avg. units/vintage \\
\midrule
$H=6$  & 4,022 & 4,022 & 304 & 1.0 & 1.51 & 0 (0.0\%) & 13.2 \\
$H=12$ & 4,034 & 4,034 & 304 & 1.0 & 1.51 & 0 (0.0\%) & 13.3 \\
$H=24$ & 4,058 & 4,058 & 304 & 1.0 & 1.51 & 0 (0.0\%) & 13.3 \\
$H=36$ & 4,086 & 4,086 & 304 & 1.0 & 1.50 & 0 (0.0\%) & 13.4 \\
\bottomrule
\end{tabular}
\end{adjustbox}
\par\medskip
\ExhibitNote{The table reports release-validation diagnostics for alternative frontier windows. Window length $H$ is measured in monthly reference-date distance from the vintage endpoint $e(\tau)$. The baseline $H=6$ window identifies the current ragged-edge completion block used in the main empirical illustration. Wider windows are diagnostic and may include older gaps, benchmark revisions, or backfill-like missingness. Reveal delays are measured in vintage steps among validated units.}
\end{table}

The diagnostic also shows that widening the frontier window changes the empirical sample only marginally in this FRED-MD implementation. This supports the use of the six-month window as a transparent current-frontier convention rather than as a tuning choice that drives the release-validation results.

\begin{figure}[H]
\centering
\begin{subfigure}{0.48\textwidth}
\centering
\includegraphics[width=\textwidth]{IA_fig_frontier_units.pdf}
\caption{Frontier units}
\end{subfigure}\hfill
\begin{subfigure}{0.48\textwidth}
\centering
\includegraphics[width=\textwidth]{IA_fig_frontier_delay.pdf}
\caption{Reveal delay}
\end{subfigure}
\caption{Frontier-window profiles}
\label{fig:ia_frontier_window_profiles}
\ExhibitNote{The figure reports how the frontier decision population and reveal-delay profile vary with the frontier window. The baseline window $H=6$ is a design choice for current ragged-edge completion. Wider windows are retained as diagnostics because they can admit older historical gaps and backfill-like cells.}
\end{figure}

\section{Method-Class Implementation}
\label{ia:procedure_classes}

The empirical menu contains fixed vintage-feasible completion rules, grouped into method classes. The sequential selector is not an additional completion rule; it is the state-to-method rule that chooses among the fixed classes using the validation archive available at the decision vintage. Each class is documented at the level needed for replication: the vintage-feasible inputs, the implementation rule, fallback behavior, and the real-time restriction.

The descriptions are intentionally class-level rather than code-level. The objective is to record which information each method may use at the decision vintage and how failures of the preferred input are handled. Detailed numerical tolerances, convergence settings, and archived outputs belong to the replication package.

\begin{table}[H]
\centering
\caption{Method classes}
\label{tab:ia_procedure_classes}
\small
\begin{adjustbox}{max width=\textwidth}
\begin{tabular}{@{}P{0.20\textwidth}P{0.25\textwidth}P{0.29\textwidth}P{0.20\textwidth}@{}}
\toprule
Method class & Vintage-feasible inputs & Implementation rule & Real-time restriction \\
\midrule
Local-history & Own-series history observed in the decision vintage & Last available same-series observation before the reference date & Uses only observations available before validation at vintage $\tau$ \\
Cross-sectional pooling & Cross-section observed at the reference date in the decision vintage & Standardized cross-sectional mean at the reference date, back-transformed to the target-series scale & Uses only contemporaneous values observed at vintage $\tau$ \\
Common-component & Vintage panel observed by the decision vintage & Two-factor low-rank reconstruction after vintage standardization and mean initialization of missing cells & Factors are estimated only from the vintage-$\tau$ panel \\
Dynamic-filtering & Own-series one-sided dynamic history observed by the decision vintage & One-sided exponentially weighted local dynamic update & Filtering is one-sided and never uses future validation values \\
\bottomrule
\end{tabular}
\end{adjustbox}
\par\medskip
\ExhibitNote{The table documents the fixed method classes used in the empirical release-validation and target-object layers. Fallback rules use only vintage-feasible same-series or panel means when the preferred input for a class is unavailable. All candidate completions are measurable with respect to the decision-vintage information set. The sequential selector maps states into these fixed method classes and is not an additional completion rule.}
\end{table}

\section{Release-Validation Diagnostics}
\label{ia:release_validation_diagnostics}

This section documents the state construction and learning diagnostics behind the release-validation results. The main text reports that the common-component method is the strongest fixed release-validation benchmark and that the feasible selector remains close to it without improving on it in local release loss. The diagnostics below explain why this happens: most frontier decisions fall in states where the common-component rule is also the state oracle, and the only state favoring dynamic filtering is comparatively small.

\begin{table}[H]
\centering
\caption{State counts}
\label{tab:ia_state_counts}
\small
\begin{adjustbox}{max width=\textwidth}
\begin{tabular}{@{}llrrrr@{}}
\toprule
Operational state & State definition & Frontier units & Validated units & Share & First initialized vintage \\
\midrule
gap0 medium & gap0; coverage medium & 2,339 & 2,339 & 58.2\% & 2000-01 \\
gap1--2 pooled & gap1--2; coverage pooled & 874 & 874 & 21.7\% & 2000-02 \\
gap0 low & gap0; coverage low & 532 & 532 & 13.2\% & 2001-06 \\
gap3--6 high & gap3--6; coverage high & 277 & 277 & 6.9\% & 2001-06 \\
\bottomrule
\end{tabular}
\end{adjustbox}
\par\medskip
\ExhibitNote{The table reports the operational state partition used by the feasible release-validation selector. States are based on vintage-available gap and coverage information, with sparse raw cells pooled by the rule documented in the data archive. All 4,022 frontier units in the six-month window ($H=6$) are release-validated in this sample. The first initialized vintage is the first vintage at which the state meets the prespecified archive threshold.}
\end{table}

The state partition is intentionally coarse. A finer partition would create more state cells but would also increase initialization and estimation costs. The table shows that the release-validation sample is concentrated in the current-gap medium-coverage state, while the high-coverage longer-gap state is much smaller.

\begin{table}[H]
\centering
\caption{State risks}
\label{tab:ia_state_risks}
\small
\begin{adjustbox}{max width=\textwidth}
\begin{tabular}{@{}lrrrrrl@{}}
\toprule
State & Records & Local-history & Pooling & Common-component & Dynamic-filtering & State oracle class \\
\midrule
gap0 medium & 2,339 & 2.447 & 1.056 & 0.932 & 1.247 & common-component \\
gap1--2 pooled & 874 & 2.021 & 1.479 & 1.384 & 1.461 & common-component \\
gap0 low & 532 & 2.886 & 2.049 & 1.532 & 2.075 & common-component \\
gap3--6 high & 277 & 1.583 & 1.790 & 1.692 & 1.309 & dynamic-filtering \\
\bottomrule
\end{tabular}
\end{adjustbox}
\par\medskip
\ExhibitNote{Entries are mean squared local release-validation losses by operational state and fixed method class. Lower values indicate better local release performance. The state oracle class is infeasible and is reported only to summarize the amount of state-dependent heterogeneity available to a feasible selector.}
\end{table}

The risk table summarizes the empirical mechanism. The common-component class is the state oracle in the three largest states. Dynamic filtering is best only in the smaller longer-gap high-coverage state. Hence the population value of state dependence is positive but limited, and a feasible selector must learn a small state-specific departure from an already strong fixed benchmark.

\begin{table}[H]
\centering
\caption{Selector diagnostics}
\label{tab:ia_selector_diagnostics}
\small
\begin{adjustbox}{max width=\textwidth}
\begin{tabular}{@{}lll@{}}
\toprule
Diagnostic & Value & Interpretation \\
\midrule
Initialization vintage & 2000-01 & First vintage satisfying the state archive threshold \\
Uninitialized target share & 3.1\% & Share of decisions using the default common-component rule \\
Oracle-static gain & 0.026 & Best fixed risk minus infeasible state-oracle risk \\
Selector-static gap & 0.027 & Feasible selector risk minus best fixed risk \\
Largest selector assignment share & 95.0\% & Largest method share assigned by the feasible selector \\
Selector assignment shares & local 0.0\%; pooling 1.1\%; common 95.0\%; dynamic 3.9\% & Assignment frequencies across decisions in the $H=6$ window \\
\bottomrule
\end{tabular}
\end{adjustbox}
\par\medskip
\ExhibitNote{The table connects the empirical release-validation result to the gain-cost decomposition. The selector uses only validation losses that arrived before each decision vintage. The common-component method is often close to the state oracle, so the estimated oracle-static gain is small and the feasible selector need not dominate the best fixed benchmark.}
\end{table}

The selector diagnostics are consistent with the state-risk evidence. The feasible rule assigns the common-component method in most decisions and deviates from it only rarely. The selector-static gap is the finite-sample cost of learning a small amount of state dependence.

\begin{figure}[H]
\centering
\includegraphics[width=0.66\textwidth]{IA_fig_selector_shares.pdf}
\caption{Selector shares}
\label{fig:ia_selector_shares}
\ExhibitNote{The figure reports assignment shares of the feasible release-validation selector by operational state. The selector chooses among fixed method classes and is not an additional completion rule. The assignment pattern illustrates why the selector remains close to the common-component benchmark when state-dependent gain is small.}
\end{figure}

\section{Target-Object Design}
\label{ia:target_object_design}

The downstream target-object exercise evaluates each completed panel through the same target equation. This section records the timing and loss object for that exercise. Its purpose is not to introduce a new forecasting model, but to show how local release-validation rankings can be compared with a downstream loss computed from a completed panel. For a fixed method class $k$, let $\widehat Y^{(k)}_\tau$ denote the panel completed at vintage $\tau$. The corresponding prediction or estimate of the target object $Z_{\tau+h}$ is
\[
\widehat Z^{(k)}_{\tau+h\mid\tau}=\Phi_{\tau,h}\!\left(\widehat Y^{(k)}_\tau,\Fcal_\tau\right),
\]
and the target loss is
\[
L^Z_\tau(k)=L^Z\!\left(\widehat Z^{(k)}_{\tau+h\mid\tau},Z_{\tau+h}\right).
\]
The map $\Phi_{\tau,h}$ is vintage-available and common across completed panels. Thus the target-object comparison changes the completion method or selection rule, not the downstream equation. This convention isolates the contribution of the completion layer. Any difference in target-object loss is attributable to the completed panel supplied to the common downstream equation, not to changing the target model across methods.

\begin{table}[H]
\centering
\caption{Target design}
\label{tab:ia_target_design}
\small
\begin{adjustbox}{max width=\textwidth}
\begin{tabular}{@{}ll@{}}
\toprule
Item & Design choice \\
\midrule
Target object & INDPRO \\
Horizon & One month \\
Forecast origins/evaluation points & 304 monthly evaluation origins \\
Target equation & Factor-augmented AR(1) for INDPRO \\
Predictors & Constant, lagged INDPRO, first principal component of the completed panel \\
Completion variation & Fixed method class or feasible release-validation selection rule \\
Target realization timing & First later vintage in which the target reference date is observed \\
\bottomrule
\end{tabular}
\end{adjustbox}
\par\medskip
\ExhibitNote{The table audits the target-object layer. The target is INDPRO and the horizon is one reference month ahead from the vintage endpoint. The same downstream equation is applied to every completed panel, so differences in target-object loss reflect the completion layer rather than changes in forecasting specification. Target reference dates are evaluated only when the corresponding INDPRO validation value appears in a later vintage.}
\end{table}

\begin{table}[H]
\centering
\caption{Target losses}
\label{tab:ia_target_losses}
\small
\begin{adjustbox}{max width=\textwidth}
\begin{tabular}{@{}lrrrrrr@{}}
\toprule
Method or rule & Local RMSE & Target MSE & Target RMSE & Target MAE & Relative target MSE & Difference vs benchmark \\
\midrule
Local-history & 1.534 & 1.812 & 1.346 & 0.663 & 1.001 & 0.002 \\
Cross-sectional pooling & 1.153 & 1.792 & 1.338 & 0.664 & 0.990 & -0.019 \\
Common-component & 1.078 & 1.810 & 1.345 & 0.665 & 1.000 & 0.000 \\
Dynamic-filtering & 1.186 & 1.795 & 1.340 & 0.663 & 0.992 & -0.015 \\
Best fixed release-validation benchmark & 1.078 & 1.810 & 1.345 & 0.665 & 1.000 & 0.000 \\
Sequential selector & 1.091 & 1.810 & 1.345 & 0.665 & 1.000 & 0.000 \\
\bottomrule
\end{tabular}
\end{adjustbox}
\par\medskip
\ExhibitNote{The table reports target-object losses for the downstream INDPRO equation. Relative target MSE is normalized to the target MSE of the best fixed release-validation benchmark, common-component. Local RMSE is reported for comparison with the release-validation layer. The sequential selector is a decision rule, not an additional completion rule, and bootstrap inference for prespecified differences is reported separately.}
\end{table}

Table~\ref{tab:ia_target_losses} reports target-object losses together with local RMSE. The comparison shows that local release accuracy and downstream target-object loss are close but not identical criteria. Cross-sectional pooling and dynamic filtering slightly improve the downstream MSE relative to the local release-validation benchmark, but the differences are small and are summarized inferentially in Section~\ref{ia:bootstrap}.

\begin{figure}[H]
\centering
\includegraphics[width=0.72\textwidth]{IA_fig_local_vs_target_loss.pdf}
\caption{Local versus target loss}
\label{fig:ia_local_vs_target_loss}
\ExhibitNote{The figure plots local release RMSE against INDPRO target-object RMSE for the fixed method classes and the feasible selector. The downstream target equation is common across completed panels. The comparison illustrates that release-validation rankings and target-object rankings need not coincide exactly.}
\end{figure}



\section{Bootstrap Inference}
\label{ia:bootstrap}

The bootstrap is an inferential layer for prespecified empirical comparisons. It is not part of the sequential selection rule and is not used to update method choices. The release-validation selector is constructed from validation losses that have arrived by the decision date; bootstrap calculations are applied only after the evaluation sample has been constructed. This separation is important. The bootstrap does not update the selection rule, choose states, select comparisons, or search for a better method. It summarizes the uncertainty in average loss differences generated by an already fixed real-time experiment.

The inferential target is limited by design. The bootstrap is used to describe sampling uncertainty in pairwise loss-difference averages, not to establish uniform dominance of one method over the whole menu. This places the target-object exercise in the tradition of predictive-ability and forecast-comparison inference, including \citet{DieboldMariano1995}, \citet{West1996}, \citet{ClarkMcCracken2001}, and \citet{GiacominiWhite2006}. The real-time nature of the exercise follows the logic of vintage-data evaluation emphasized by \citet{CroushoreStark2001}. The dependent-data resampling device follows the block-bootstrap tradition of \citet{Kunsch1989}, \citet{PolitisRomano1994}, and \citet{Lahiri2003}.

There are two bootstrap problems. The first concerns local release-validation losses. The observations are frontier decisions $u=(i,t,\tau)$, and many frontier decisions share the same decision vintage, release environment, macro shock, or panel state. The relevant resampling unit is therefore a calendar block of decision vintages, with all frontier decisions inside the sampled block kept together. The second concerns target-object losses. The observations are forecast or nowcast origins $\tau$, and the object is a time-ordered sequence of loss differences. This second problem is closest to \citet{GoncalvesMcCrackenYao2025}, who develop a block-bootstrap approach to out-of-sample predictive-ability inference with real-time data and regular revisions. Our use of their logic is conservative: we resample blocks of forecast origins after the real-time evaluation sequence has been fixed, and we use the resulting distribution only to summarize uncertainty in prespecified target-object loss differences.

All bootstrap comparisons below are pairwise and prespecified. The reported $p$-values are descriptive two-sided bootstrap summaries, not post-selection evidence from searching over many methods. The baseline block length is six monthly origins. This choice is fixed before inspecting significance patterns, is not selected to optimize any reported comparison, and is intended to preserve short-run dependence in monthly loss differences. The same number of replications, $B^\ast=999$, is used for all reported comparisons. In this sense, the bootstrap layer is closest to conditional predictive-ability inference: the realized selector assignments and the realized evaluation path are treated as fixed when forming the bootstrap distribution.

\subsection{Bootstrap statistic and interpretation}
\label{ia:bootstrap_scope}

This subsection records the statistic and interpretation behind the bootstrap calculations. They are used for prespecified empirical comparisons between completion methods or rules, not for the sequential learning algorithm itself. For a target-object comparison between two methods or rules $a$ and $b$, let
\[
    d_\tau(a,b)=L^Z_\tau(a)-L^Z_\tau(b),\qquad \tau\in\Tcal_Z,
\]
where $\Tcal_Z$ is the set of target-object evaluation origins. Write $T_Z=|\Tcal_Z|$ and
\[
    \bar d(a,b)=T_Z^{-1}\sum_{\tau\in\Tcal_Z}d_\tau(a,b).
\]
Let
\[
    \bar\mu_d=T_Z^{-1}\sum_{\tau\in\Tcal_Z}\E[d_\tau(a,b)].
\]
The expectation is taken with respect to the underlying, potentially heterogeneous real-time data-generating process, so the mean loss difference may vary across evaluation origins. The first-order approximation concerns
\[
    \sqrt{T_Z}\{\bar d(a,b)-\bar\mu_d\}.
\]
The bootstrap standard errors and intervals are computed from the conditional distribution of the centered statistic
\[
    \bar d^\ast(a,b)-\bar d(a,b),
\]
where $\bar d^\ast(a,b)$ is obtained by the moving-block bootstrap: with block length $\ell_B$, the overlapping blocks of $\ell_B$ consecutive evaluation origins, with starting points $s=1,\ldots,T_Z-\ell_B+1$, are drawn independently and uniformly with replacement, concatenated, and truncated to length $T_Z$, and $\bar d^\ast(a,b)$ is the average of the resampled loss differences. Validity in the present setting does not require stationarity, but it does require weak dependence, uniform moment bounds, and a non-degenerate limit. The following high-level condition is maintained for every prespecified comparison.

\begin{iaassumption}
\label{ia:ass-bootstrap}
Let $W_\tau$ denote the origin-level observation used in a given comparison: $W_\tau=d_\tau(a,b)$ for target-object comparisons, and $W_\tau=(D_\tau(a,b),M_\tau)'$ for release-validation comparisons, with the sample size $T_Z$ read as the number of decision vintages in the release layer.
\begin{enumerate}[label=(\alph*), leftmargin=*]
    \item $\{W_\tau\}$ is a possibly heterogeneous, weakly dependent array: strong mixing, or near-epoch dependent on a mixing array, with mixing or dependence coefficients of the sizes and uniformly bounded moments of order $2+\varepsilon$, $\varepsilon>0$, required by moving-block theory for dependent heterogeneous arrays.
    \item The long-run covariance matrix of the normalized centered partial sums of $\{W_\tau\}$ converges to a finite limit that is positive along the contrast of interest.
    \item In the release layer, $\liminf_{T_Z\to\infty}\,T_Z^{-1}\sum_{\tau}\E[M_\tau]>0$.
    \item The block length satisfies $\ell_B\to\infty$ and $\ell_B/T_Z\to0$.
\end{enumerate}
\end{iaassumption}

\begin{iaproposition}
\label{ia:prop-bootstrap}
Suppose Assumption~\ref{ia:ass-bootstrap} holds. Then, as $T_Z\to\infty$:
\begin{enumerate}[label=(\roman*), leftmargin=*]
    \item $\sqrt{T_Z}\{\bar d(a,b)-\bar\mu_d\}$ is asymptotically normal, and the moving-block bootstrap is consistent:
    \[
        \sup_{x\in\mathbb R}\Big|\,\Prb^\ast\big[\sqrt{T_Z}\{\bar d^\ast(a,b)-\bar d(a,b)\}\le x\big]
        -\Prb\big[\sqrt{T_Z}\{\bar d(a,b)-\bar\mu_d\}\le x\big]\Big|\to_p 0 ,
    \]
    where $\Prb^\ast$ denotes the bootstrap distribution conditional on the sample.
    \item The same conclusions hold for the release-validation ratio
    \[
      \bar\delta^Y(a,b)=\frac{\sum_\tau D_\tau(a,b)}{\sum_\tau M_\tau},
      \qquad
      \bar\mu_\delta=\frac{\sum_\tau\E[D_\tau(a,b)]}{\sum_\tau\E[M_\tau]},
    \]
    and its bootstrap analogue, by the delta method applied to the vintage-level mean vector.
\end{enumerate}
\end{iaproposition}

Proposition~\ref{ia:prop-bootstrap} assembles known results and claims no new bootstrap theory. Part (i) follows from central limit theory and moving-block consistency for dependent \emph{heterogeneous} arrays: \citet{GoncalvesWhite2002} establish bootstrap consistency for means of near-epoch-dependent, heterogeneously distributed data, and \citet{Fitzenberger1998} covers moving-block inference under heteroskedasticity and autocorrelation; \citet{Kunsch1989} and \citet{Lahiri2003} provide the stationary background. The heterogeneous-array results, rather than stationary theory alone, are the relevant justification here because the evaluation window spans the 2008--2009 and 2020 episodes, so the loss-difference array is best treated as non-identically distributed. Part (ii) follows by applying part (i) to the vector $W_\tau=(D_\tau,M_\tau)'$ and the map $(x,y)\mapsto x/y$, which is continuously differentiable at the limit of the vintage-level means because the denominator is bounded away from zero by part (c); the joint block resampling of $(D_\tau,M_\tau)$ is what makes this delta step valid, since it preserves the within-vintage dependence between loss sums and validation counts.

Interval and $p$-value constructions are fixed as follows. Confidence intervals are equal-tailed percentile intervals from the bootstrap distribution of the average; descriptive two-sided $p$-values are symmetric and computed from the centered bootstrap distribution, $p=(B^\ast)^{-1}\sum_{b=1}^{B^\ast}\1\{|\bar d^\ast_b-\bar d|\ge|\bar d|\}$. Because the two constructions answer slightly different questions, an interval endpoint near zero need not correspond to a small $p$-value, especially when the bootstrap distribution is asymmetric or highly concentrated. In the empirical implementation, the block length is fixed at six monthly origins to match the short-run data-release cycle and is used only as a finite-sample uncertainty summary. The resulting intervals are not a new asymptotic result about optimal block-length choice or about the full sequential learning rule.

For release-validation comparisons, let $V^Y_\tau$ be the set of release-validated frontier decisions made at decision vintage $\tau$. For a comparison between rules or method classes $a$ and $b$, define
\[
    D_\tau(a,b)=\sum_{u\in V^Y_\tau}\delta^Y_u(a,b),
    \qquad
    M_\tau=|V^Y_\tau|.
\]
Then the release-validation average can be written as the ratio
\[
    \bar\delta^Y(a,b)
    =
    \frac{\sum_\tau D_\tau(a,b)}{\sum_\tau M_\tau}.
\]
In each bootstrap replication, blocks of decision vintages are resampled and the bootstrap analogue is
\[
    \bar\delta^{Y,\ast}(a,b)
    =
    \frac{\sum_{\tau\in\mathcal B^\ast}D_\tau(a,b)}
    {\sum_{\tau\in\mathcal B^\ast}M_\tau},
\]
where $\mathcal B^\ast$ denotes the resampled multiset of decision vintages induced by the sampled blocks. This ratio form keeps all frontier decisions within a sampled vintage block together and preserves the empirical weighting induced by unequal numbers of frontier decisions across vintages. Because $M_\tau$ varies across decision vintages, the release-validation bootstrap resamples the vintage-level vector $(D_\tau(a,b),M_\tau)$ jointly rather than resampling normalized vintage averages. Assumption~\ref{ia:ass-bootstrap}(c) and Proposition~\ref{ia:prop-bootstrap}(ii) record the conditions under which this ratio scheme is asymptotically valid.

The distinction between conditional and unconditional inference is essential. Conditional on the realized decision path, the bootstrap describes uncertainty in the loss differences delivered by that path. It does not account for the additional variation that would arise if the entire sequential selection experiment were repeated from scratch and the selector followed a different path. One further distinction sharpens the scope. For comparisons between fixed method classes, no data-driven selection is involved at all: under delayed full validation, every class is scored on every validated unit, so the loss-difference sequence is a fixed function of the data and raises no post-selection concerns. Only the rows involving the feasible selector depend on the realized learning path, and for those the conditional interpretation is the operative one. The resulting intervals are therefore interpreted as finite-sample uncertainty summaries for prespecified loss differences, and for gain-style diagnostics when the relevant benchmark assignments are treated as fixed or sufficiently separated, rather than as a new asymptotic theory for the selected method or for the full sequential learning path. For gain-style quantities involving minima, the bootstrap output is reported as a diagnostic uncertainty summary rather than as formal inference for a nonsmooth argmin functional. The selected method identity is an argmin object and is not bootstrapped here; its asymptotic behavior is governed by the margin conditions discussed in the main text. This limitation matches the role of the empirical exercise: the bootstrap supports cautious comparisons among prespecified rules and methods, while the theoretical learning results in the main text describe the broader gain-cost logic of sequential selection.

\subsection{Release-validation differences}
\label{ia:release_bootstrap}

For two prespecified rules or method classes $a$ and $b$, define
\[
\delta^Y_u(a,b)=\ell^Y_u(a)-\ell^Y_u(b),
\qquad
\bar\delta^Y(a,b)=\frac{1}{|\Vcal_Y|}\sum_{u\in\Vcal_Y}\delta^Y_u(a,b),
\]
where $\Vcal_Y$ is the set of release-validated frontier decisions. Negative values favor $a$. The comparison is local: it concerns release-validation loss, not downstream target-object loss.

The release-validation bootstrap resamples blocks of consecutive decision vintages. For each sampled block, all release-validated frontier decisions whose decision vintages fall in the block are retained. This preserves the dependence induced by common vintage information, release-calendar clustering, repeated completions of related cells, and cross-sectional macro shocks. The feasible selector is not re-estimated inside each bootstrap sample. It is treated as the real-time rule produced by the original sequential selection experiment, and the bootstrap summarizes uncertainty in the resulting loss-difference sequence. Thus the release-validation bootstrap is conditional on the realized training and decision path of the selector. It does not approximate the sampling distribution that would arise from rerunning the full sequential selection experiment on new real-time panels.

\begin{table}[H]
\centering
\caption{Release bootstrap}
\label{tab:ia_bootstrap_release}
\small
\begin{adjustbox}{max width=\textwidth}
\begin{tabular}{@{}lrrrrrr@{}}
\toprule
Comparison & Mean diff. & Bootstrap SE & 95\% CI & $p$-value & Block length & Replications \\
\midrule
Selector vs best fixed release-validation benchmark & 0.0273 & 0.0505 & $[-0.0532,\ 0.1449]$ & 0.711 & 6 & 999 \\
State oracle vs best fixed release-validation benchmark & -0.0263 & 0.0228 & $[-0.0737,\ 0.0107]$ & 0.238 & 6 & 999 \\
\bottomrule
\end{tabular}
\end{adjustbox}
\par\medskip
\ExhibitNote{Negative values favor the first rule in the comparison. Loss differences use local release-validation losses on frontier units in the six-month ($H=6$) window. The bootstrap resamples calendar blocks of decision vintages and preserves all frontier units inside sampled blocks. The selector comparison concerns the feasible real-time rule, conditional on the realized decision path. The state-oracle comparison is diagnostic and infeasible: it summarizes uncertainty around the estimated amount of state-dependent heterogeneity, not the performance of an implementable rule. Confidence intervals are equal-tailed percentile intervals; descriptive two-sided $p$-values are symmetric and computed from the centered bootstrap distribution (Section~\ref{ia:bootstrap_scope}). Confidence-interval endpoints are rounded to four decimal places.}
\end{table}

Table~\ref{tab:ia_bootstrap_release} reinforces the main release-validation interpretation. The feasible selector is close to the best fixed benchmark, but the estimated difference is small relative to bootstrap uncertainty. The infeasible state oracle improves on the best fixed benchmark on average, but its interval also includes zero. This is the empirical counterpart of the gain-cost logic: the state-dependent gain exists, but it is small in this FRED-MD frontier window.

\subsection{Target-object differences}
\label{ia:target_bootstrap}

For two prespecified methods or rules $k$ and $m$, define
\[
    d_\tau(k,m)=L^Z_\tau(k)-L^Z_\tau(m),\qquad \tau\in\Tcal_Z,
\]
where $\Tcal_Z$ is the set of target-object evaluation origins. Negative values favor method $k$. The average difference is
\[
    \bar d(k,m)=\frac{1}{|\Tcal_Z|}\sum_{\tau\in\Tcal_Z}d_\tau(k,m).
\]
The target-object comparison is the closest analogue to real-time forecast-comparison inference. The completed panel, the target equation, and the target realization are all indexed by the real-time evaluation origin. \citet{GoncalvesMcCrackenYao2025} show how block bootstrap methods can be used for predictive-ability inference when forecasts are constructed from real-time vintage data subject to regular revisions. Our implementation follows the same organizing principle: once the sequence of real-time target-object loss differences has been fixed, we resample blocks of consecutive evaluation origins to preserve serial dependence, vintage dependence, and revision-related dependence in the loss-difference sequence.

The target-object bootstrap is not a bootstrap of the entire sequential selection experiment. In particular, it does not rerun the release-validation selector inside each resample. This is intentional. The object reported in Table~\ref{tab:ia_bootstrap_target} is the uncertainty around prespecified average target-object loss differences for the completed panels generated by the original real-time exercise. As in the release-validation layer, the approximation is conditional on the realized real-time path: the completed panels, target equation, selector decisions, forecast origins, and target realizations used to form the loss-difference sequence are fixed before resampling begins.

\begin{table}[H]
\centering
\caption{Target bootstrap}
\label{tab:ia_bootstrap_target}
\small
\begin{adjustbox}{max width=\textwidth}
\begin{tabular}{@{}lrrrrrr@{}}
\toprule
Comparison & Mean diff. & Bootstrap SE & 95\% CI & $p$-value & Block length & Replications \\
\midrule
Pooling vs common-component & -0.0187 & 0.0178 & $[-0.0599,\ 0.0033]$ & 0.162 & 6 & 999 \\
Dynamic-filtering vs common-component & -0.0152 & 0.0144 & $[-0.0478,\ 0.0020]$ & 0.116 & 6 & 999 \\
Selector vs common-component & 0.0002 & 0.0002 & $[0.0000,\ 0.0005]$ & 0.330 & 6 & 999 \\
Local-history vs common-component & 0.0017 & 0.0049 & $[-0.0051,\ 0.0131]$ & 0.785 & 6 & 999 \\
\bottomrule
\end{tabular}
\end{adjustbox}
\par\medskip
\ExhibitNote{Negative values favor the first method or rule. Loss differences are computed on downstream INDPRO target-object squared loss, not on local release loss. The moving block bootstrap resamples consecutive evaluation origins and is used only to summarize uncertainty for prespecified comparisons against the common-component benchmark. Confidence intervals are equal-tailed percentile intervals; descriptive two-sided $p$-values are symmetric, computed from the centered bootstrap distribution (Section~\ref{ia:bootstrap_scope}), and not adjusted for searching over multiple comparisons. The selector coincides with the common-component completion at almost all origins, so its loss-difference sequence is zero for most $\tau$; the corresponding bootstrap distribution is highly concentrated near zero, and that row should be read as documenting near-equivalence rather than as a precise interval. Confidence-interval endpoints are rounded to four decimal places.}
\end{table}

Table~\ref{tab:ia_bootstrap_target} supports a cautious reading of the downstream exercise. Pooling and dynamic filtering have lower average target-object loss than the common-component benchmark, but the bootstrap intervals include zero. The selector is essentially indistinguishable from the common-component method on the target-object criterion. The table therefore documents the validation-target distinction without turning the empirical illustration into a claim of target-object dominance.

\begin{algorithm}[H]
\caption{Target-object block bootstrap}
\label{alg:ia_target_bootstrap}
\begin{algorithmic}[1]
\State Fix a prespecified pair $(k,m)$ and compute $d_\tau(k,m)$ over target-object evaluation origins.
\State Choose a block length $\ell_B$ and number of bootstrap replications $B^\ast$ before inspecting significance patterns.
\For{$b=1,\ldots,B^\ast$}
  \State Draw blocks of consecutive evaluation origins until the bootstrap sequence has length at least $|\Tcal_Z|$.
  \State Truncate to length $|\Tcal_Z|$ and compute the bootstrap average loss difference.
\EndFor
\State Report the point estimate, bootstrap standard error, confidence interval, and descriptive two-sided $p$-value.
\end{algorithmic}
\end{algorithm}

\section{Additional Monte Carlo Output}
\label{ia:additional_mc}

This section collects simulation material that is useful for auditability but too detailed for the main text. The main paper reports the core Monte Carlo evidence and discusses the gain-cost mechanism. The appendix records the calibration, bridge diagnostics, scaling exercises, application-size regimes, and bridge-strength diagnostics so that the simulation claims in the main text can be checked without interrupting the main narrative.

\begin{table}[H]
\centering
\caption{Monte Carlo calibration}
\label{tab:ia_mc_calibration}
\small
\begin{adjustbox}{max width=\textwidth}
\begin{tabular}{@{}p{0.27\textwidth}p{0.66\textwidth}@{}}
\toprule
Design component & Calibration \\
\midrule
Baseline size & $R=1000$, $N_0=60$, $T_0=250$. \\
Panel DGP & $Y_{it}=\lambda_i f_t+u_{it}$, $u_{it}=\rho_i u_{i,t-1}+\sigma_i\varepsilon_{it}$; heterogeneity in $(\lambda_i,\rho_i,\sigma_i)$ generates state-dependent rankings. \\
States & $G_u\in\{F,P,L\}$: factor-dominant, persistence-dominant, and noisy/local cells; $G_u$ is observed at the decision date. \\
Method menu & $\Kcal=\{\text{Local},\text{Pool},\text{Factor}\}$: own-lag persistence, lagged cross-sectional pooling, and one-step factor forecast. \\
Feasibility & $\widehat Y^{(k)}_{it|\tau}=q_k(u;\Fcal_\tau)$; the validation value $Y_{it}$ enters only when validation loss is computed. \\
Initialization & Full validation: $m=2$ validated records per state, equivalently per state-method cell because all methods are scored on each validated unit; selected-arm feedback: $m=2$ observed records per state-method cell. Sparse-feedback design: $m_{S2}=20$. \\
Delays & $(\bar D,D_{\max})=(2,8)$ in S0, S1, S4, S5, and S6; $(22,70)$ in S2; $(3,12)$ in S3. \\
Feedback & Full validation: $\Kcal^V=\Kcal$; selected-arm feedback in S6: $\Kcal^V=\{A_u\}$ with exploration probability $\epsilon=0.12$. \\
S4 transport & $p^V(H)\ne p^T(H)$; correction reweights by observed subtype $H$. \\
S5 target object & Local-trained rule is learned from local validation and evaluated on target-object loss; the target-trained rule is infeasible. \\
Scaling grids & $T\in\{50,100,250,500,1000,2000\}$ at $N=60$; $N\in\{30,60,120,240,500,1000\}$ at $T=250$. \\
Application regimes & Stylized macro, finance, micro, and completed-panel dimensions mirror the cases in the main text. \\
\bottomrule
\end{tabular}
\end{adjustbox}
\par\medskip
\ExhibitNote{The table records the notation and calibration choices needed to reproduce the Monte Carlo design. The basic panel equation and method menu are common across scenarios. Scenario-specific settings change validation timing, feedback observability, separation of state rankings, or the validation-target relation.}
\end{table}

\begin{table}[H]
\centering
\caption{Bridge diagnostics}
\label{tab:ia_mc_bridge}
\scriptsize
\begin{adjustbox}{max width=\textwidth}
\begin{tabular}{@{}llrrrrr@{}}
\toprule
Scenario & Metric & Mean & MCSE & Median & p10 & p90 \\
\midrule
S3 & Cell/block rank agreement & 1.000 & 0.0000 & 1.000 & 1.000 & 1.000 \\
S4 & Unadjusted rule target risk & 1.426 & 0.004 & 1.435 & 1.244 & 1.595 \\
S4 & Transport-corrected rule target risk & 1.162 & 0.004 & 1.161 & 1.019 & 1.304 \\
S4 & Transport risk gain & 0.264 & 0.004 & 0.266 & 0.104 & 0.417 \\
S4 & Rank agreement: unadjusted vs target & 0.630 & 0.005 & 0.667 & 0.333 & 0.667 \\
S4 & Rank agreement: corrected vs target & 0.867 & 0.005 & 1.000 & 0.667 & 1.000 \\
S5 & Local-trained rule target risk & 0.706 & 0.003 & 0.699 & 0.582 & 0.837 \\
S5 & Infeasible target-trained rule risk & 0.304 & 0.001 & 0.301 & 0.258 & 0.352 \\
S5 & Target-object deterioration & 0.402 & 0.002 & 0.395 & 0.320 & 0.494 \\
S5 & Local-target rank agreement & 0.341 & 0.002 & 0.333 & 0.333 & 0.333 \\
S5 & Local-target divergence indicator & 0.659 & 0.002 & 0.667 & 0.667 & 0.667 \\
\bottomrule
\end{tabular}
\end{adjustbox}
\par\medskip
\ExhibitNote{S3 is an aligned aggregate-validation design. S4 concerns validation-population shift and transport correction. S5 concerns target-object evaluation. Local-trained rule means learned from local validation and evaluated on target-object loss. The target-trained rule is infeasible and is used only as a benchmark. Monte Carlo standard errors are reported in the adjacent MCSE column.}
\end{table}

The bridge diagnostics separate aligned validation from validation-target mismatch. S3 confirms that changing the loss-bearing object from cells to blocks need not create a problem when the block loss is also the target criterion. S4 and S5 show the opposite case: when validation and target populations or criteria differ, the local validation ranking can be misleading unless an explicit bridge or transport correction is available.

\begin{landscape}
\begin{table}[p]
\centering
\caption{$T$-scaling}
\label{tab:ia_mc_T_scaling}
\scriptsize
\begin{adjustbox}{max width=\linewidth}
\begin{tabular}{@{}lllrrrrrr@{}}
\toprule
Scenario & Metric & Stat. & $T=50$ & $T=100$ & $T=250$ & $T=500$ & $T=1000$ & $T=2000$ \\
\midrule
S1 & $G$ & Mean & 0.292 & 0.301 & 0.306 & 0.305 & 0.309 & 0.310 \\
S1 & $G$ & MCSE & 0.002 & 0.002 & 0.002 & 0.002 & 0.002 & 0.002 \\
S1 & Regret & Mean & 0.029 & 0.014 & 0.005 & 0.002 & 0.001 & 0.001 \\
S1 & Regret & MCSE & 0.0006 & 0.0003 & 0.0001 & $<0.001$ & $<0.001$ & $<0.001$ \\
S1 & Recovered & Mean & 0.893 & 0.952 & 0.982 & 0.991 & 0.996 & 0.998 \\
S1 & Recovered & MCSE & 0.002 & 0.001 & 0.0004 & 0.0002 & 0.0001 & $<0.001$ \\
S1 & Init. share & Mean & 0.025 & 0.012 & 0.005 & 0.002 & 0.001 & 0.001 \\
S1 & Archive size & Mean & 734.9 & 1554.3 & 4015.4 & 8079.6 & 16036.3 & 32408.7 \\
S1 & Sparsity & Mean & 0.187 & 0.173 & 0.170 & 0.174 & 0.187 & 0.180 \\
S1 & Oracle match & Mean & 0.945 & 0.973 & 0.990 & 0.995 & 0.998 & 0.999 \\
\midrule
S2 & $G$ & Mean & 0.238 & 0.248 & 0.254 & 0.260 & 0.260 & 0.265 \\
S2 & Regret & Mean & 0.082 & 0.038 & 0.017 & 0.008 & 0.004 & 0.002 \\
S2 & Recovered & Mean & 0.609 & 0.831 & 0.928 & 0.969 & 0.986 & 0.994 \\
S2 & Init. share & Mean & 0.141 & 0.069 & 0.027 & 0.014 & 0.007 & 0.003 \\
S2 & Archive size & Mean & 87.1 & 226.4 & 676.2 & 1459.3 & 2952.8 & 5967.0 \\
S2 & Sparsity & Mean & 0.807 & 0.819 & 0.818 & 0.813 & 0.816 & 0.814 \\
S2 & Oracle match & Mean & 0.773 & 0.823 & 0.886 & 0.939 & 0.972 & 0.990 \\
\midrule
S6 & $G$ & Mean & 0.292 & 0.300 & 0.305 & 0.307 & 0.310 & 0.311 \\
S6 & Regret & Mean & 0.149 & 0.126 & 0.101 & 0.092 & 0.083 & 0.078 \\
S6 & Recovered & Mean & 0.460 & 0.564 & 0.657 & 0.690 & 0.726 & 0.743 \\
S6 & Init. share & Mean & 0.040 & 0.020 & 0.008 & 0.004 & 0.002 & 0.001 \\
S6 & Archive size & Mean & 39.5 & 71.5 & 168.6 & 330.0 & 651.3 & 1300.3 \\
S6 & Sparsity & Mean & 0.527 & 0.465 & 0.403 & 0.375 & 0.348 & 0.333 \\
S6 & Oracle match & Mean & 0.716 & 0.750 & 0.819 & 0.846 & 0.880 & 0.898 \\
\bottomrule
\end{tabular}
\end{adjustbox}
\par\medskip
\ExhibitNote{All entries use 1,000 replications. The table varies $T$ at $N=60$. For compactness, the table reports means for the main diagnostics and reports MCSE rows only for the headline S1 quantities. Recovered is the recovered gain share; Regret is relative to the state-dependent oracle; archive size is the terminal number of arrived validation records. Increasing $T$ expands the validation archive over decision time and is especially important when validation is delayed or selected-arm feedback makes observed state-method losses sparse.}
\end{table}
\end{landscape}

\begin{landscape}
\begin{table}[p]
\centering
\caption{$N$-scaling}
\label{tab:ia_mc_N_scaling}
\scriptsize
\begin{adjustbox}{max width=\linewidth}
\begin{tabular}{@{}lllrrrrrr@{}}
\toprule
Scenario & Metric & Stat. & $N=30$ & $N=60$ & $N=120$ & $N=240$ & $N=500$ & $N=1000$ \\
\midrule
S1 & $G$ & Mean & 0.302 & 0.307 & 0.305 & 0.306 & 0.306 & 0.307 \\
S1 & Regret & Mean & 0.007 & 0.006 & 0.004 & 0.004 & 0.004 & 0.003 \\
S1 & Recovered & Mean & 0.975 & 0.981 & 0.985 & 0.987 & 0.988 & 0.989 \\
S1 & Init. share & Mean & 0.006 & 0.005 & 0.004 & 0.004 & 0.004 & 0.004 \\
S1 & Archive size & Mean & 1829.8 & 4008.2 & 8489.9 & 17765.3 & 38224.5 & 77969.4 \\
S1 & Sparsity & Mean & 0.235 & 0.170 & 0.128 & 0.092 & 0.063 & 0.045 \\
S1 & Oracle match & Mean & 0.985 & 0.989 & 0.993 & 0.993 & 0.994 & 0.995 \\
\midrule
S2 & $G$ & Mean & 0.240 & 0.254 & 0.263 & 0.269 & 0.268 & 0.271 \\
S2 & Regret & Mean & 0.021 & 0.016 & 0.013 & 0.012 & 0.011 & 0.010 \\
S2 & Recovered & Mean & 0.905 & 0.931 & 0.948 & 0.955 & 0.958 & 0.962 \\
S2 & Init. share & Mean & 0.039 & 0.027 & 0.020 & 0.014 & 0.011 & 0.008 \\
S2 & Archive size & Mean & 395.6 & 692.7 & 1369.1 & 2723.5 & 5604.5 & 11326.2 \\
S2 & Sparsity & Mean & 0.776 & 0.813 & 0.820 & 0.826 & 0.828 & 0.827 \\
S2 & Oracle match & Mean & 0.881 & 0.895 & 0.909 & 0.903 & 0.901 & 0.911 \\
\midrule
S3 & $G$ & Mean & 0.114 & 0.114 & 0.115 & 0.115 & 0.115 & 0.115 \\
S3 & Regret & Mean & 0.001 & 0.001 & 0.001 & 0.001 & 0.001 & 0.001 \\
S3 & Recovered & Mean & 0.990 & 0.992 & 0.994 & 0.994 & 0.994 & 0.994 \\
S3 & Init. share & Mean & 0.006 & 0.005 & 0.005 & 0.004 & 0.004 & 0.004 \\
S3 & Archive size & Mean & 1786.9 & 3972.3 & 8525.4 & 17658.3 & 38042.4 & 77709.6 \\
S3 & Sparsity & Mean & 0.249 & 0.179 & 0.123 & 0.097 & 0.065 & 0.045 \\
S3 & Oracle match & Mean & 0.993 & 0.994 & 0.995 & 0.995 & 0.995 & 0.996 \\
\bottomrule
\end{tabular}
\end{adjustbox}
\par\medskip
\ExhibitNote{All entries use 1,000 replications. The table varies $N$ at $T=250$. For compactness, the table reports means for the main diagnostics. Recovered is the recovered gain share; Regret is relative to the state-dependent oracle; archive size is the terminal number of arrived validation records. Increasing $N$ expands the cross-sectional dimension and stabilizes state-method feedback counts, but does not replace a long decision horizon.}
\end{table}
\end{landscape}

\begin{landscape}
\begin{table}[p]
\centering
\caption{Application regimes}
\label{tab:ia_mc_application_regimes}
\scriptsize
\begin{adjustbox}{max width=\linewidth}
\begin{tabular}{@{}llrrrrrr@{}}
\toprule
Application & Scenario & $N$ & $T$ & Recovered & Regret & Archive size & Bridge/target \\
\midrule
Macro & S1 & 60 & 250 & 0.982 & 0.005 & 3973.9 & -- \\
Macro & S2 & 60 & 250 & 0.935 & 0.016 & 696.8 & -- \\
Macro & S1 & 60 & 500 & 0.991 & 0.003 & 8061.8 & -- \\
Macro & S2 & 60 & 500 & 0.969 & 0.008 & 1393.7 & -- \\
Macro & S1 & 60 & 1000 & 0.996 & 0.001 & 16185.6 & -- \\
Macro & S2 & 60 & 1000 & 0.986 & 0.004 & 3061.3 & -- \\
Macro & S1 & 120 & 250 & 0.986 & 0.004 & 8506.7 & -- \\
Macro & S2 & 120 & 250 & 0.943 & 0.014 & 1343.0 & -- \\
Macro & S1 & 120 & 500 & 0.993 & 0.002 & 17265.5 & -- \\
Macro & S2 & 120 & 500 & 0.975 & 0.007 & 2844.9 & -- \\
Finance & S3 & 250 & 1000 & 0.999 & 0.000 & 75202.3 & -- \\
Finance & S6 & 250 & 1000 & 0.733 & 0.081 & 3016.0 & -- \\
Finance & S3 & 500 & 2000 & 0.999 & 0.000 & 310447.4 & -- \\
Finance & S6 & 500 & 2000 & 0.748 & 0.077 & 12436.4 & -- \\
Micro & S4 & 500 & 20 & 0.567 & 0.122 & 534.7 & 0.216 \\
Micro & S4 & 500 & 100 & 0.920 & 0.025 & 3303.3 & 0.295 \\
Micro & S4 & 2000 & 100 & 0.947 & 0.017 & 13154.3 & 0.306 \\
Completed & S5 & 60 & 250 & -- & 0.407 & 4016.3 & 0.403 \\
Completed & S5 & 120 & 500 & -- & 0.397 & 17225.8 & 0.395 \\
Completed & S5 & 240 & 500 & -- & 0.398 & 35822.3 & 0.396 \\
\bottomrule
\end{tabular}
\end{adjustbox}
\par\medskip
\ExhibitNote{The table reports representative application-size regimes. These are scale calibrations for the Monte Carlo designs, not additional empirical applications. Macro regimes use delayed full-validation scenarios S1 and S2. Finance regimes use aligned aggregate validation (S3) and selected-arm feedback (S6). Micro regimes use S4; the bridge/target column reports the transport diagnostic. Completed-panel regimes use S5; the bridge/target column reports target-object deterioration. The full simulation archive contains the complete grid and MCSE columns.}
\end{table}
\end{landscape}

\begin{table}[H]
\centering
\caption{Bridge strength}
\label{tab:ia_mc_bridge_strength}
\small
\begin{adjustbox}{max width=\textwidth}
\begin{tabular}{@{}llrrrrr@{}}
\toprule
Scenario & Strength & Diagnostic & Mean & MCSE & Median & p90 \\
\midrule
S4 & 0.00 & Transport risk gain & -0.000 & $<0.001$ & 0.000 & 0.000 \\
S4 & 0.50 & Transport risk gain & -0.000 & 0.0002 & -0.000 & 0.005 \\
S4 & 1.00 & Transport risk gain & 0.256 & 0.004 & 0.261 & 0.409 \\
S4 & 1.50 & Transport risk gain & 0.335 & 0.006 & 0.320 & 0.593 \\
S4 & 2.00 & Transport risk gain & 0.389 & 0.008 & 0.369 & 0.710 \\
S5 & 0.00 & Target-object deterioration & 0.190 & 0.002 & 0.184 & 0.258 \\
S5 & 0.50 & Target-object deterioration & 0.298 & 0.002 & 0.297 & 0.374 \\
S5 & 1.00 & Target-object deterioration & 0.402 & 0.002 & 0.395 & 0.494 \\
S5 & 1.50 & Target-object deterioration & 0.503 & 0.003 & 0.502 & 0.606 \\
S5 & 2.00 & Target-object deterioration & 0.603 & 0.003 & 0.597 & 0.730 \\
\bottomrule
\end{tabular}
\end{adjustbox}
\par\medskip
\ExhibitNote{S4 varies validation-target mismatch strength and reports the risk improvement from transport correction. S5 varies target-object divergence strength and reports the excess target-object risk of the local-trained rule relative to the infeasible target-trained benchmark. All entries use 1,000 replications.}
\end{table}

The bridge-strength table makes the identification message explicit. In S4, transport correction matters only when validation and target compositions diverge. In S5, target-object deterioration grows with the strength of the local-target divergence. These are not learning failures: they are consequences of evaluating a rule under a target criterion that is not aligned with the validation criterion used to learn it.

\bibliographystyle{elsarticle-harv}
\bibliography{References_BfP}

\end{document}
